\chapter{The social contract}
\label{ch:welfare}

\section{The commitment}

Cuba keeps a universal welfare state at Western European standard --- health,
education, pensions, disability, child support, income support --- and the rest
of Part III is arranged around being able to afford it.

That is the book's one non-negotiable, and this chapter is where it is argued
rather than asserted. Four reasons, in descending order of how much weight they
carry.

\emph{It is the thing worth keeping.} Chapter~\ref{ch:ledgerch} audited
sixty-seven years and found one unambiguous asset. The health outcomes are
measured, not claimed; the education results were tested by an outside body and
beat the region by two standard deviations, with the poorest quartile above the
regional average. Nothing else in the ledger is in that class.

\emph{It is what makes the rest politically possible.} A population asked to
accept markets, foreign owners, price liberalisation and privatisation will
accept them if the floor holds. It will not if the floor goes, and the evidence
is not theoretical: the post-communist transitions that withdrew social
protection early had their economic reforms reversed by their own electorates
within a decade \citep{roland2000}, which is the most expensive way to discover that consent is an
input. Cuba is being asked to run a transition with a population that has been
told for sixty-seven years that this is what the sacrifice bought. Take it away
in year two and there is no year five.

\emph{It is the protection for the people the transition will hurt.}
Chapter~\ref{ch:premises} identified who is exposed: households without
remittances, which are disproportionately Black; the eastern provinces; the old;
the rural. Every liberalising measure in this book advantages the household with
capital and a relative abroad. The universal floor is the instrument that stops
that from becoming a permanent stratification, and it is the only instrument
that works without an administrative capacity Cuba does not have.

\emph{It is an input, not only a transfer.} The education system is what
Chapter~\ref{ch:talent} and Chapter~\ref{ch:life} propose to monetise. It cannot
be defunded for a decade and restarted, because the teachers leave and the
cohorts pass through. Spending on it is not consumption of the growth model's
proceeds; it is the growth model's raw material.

\section{What changes is who pays}

Here is the chapter's central move, and the sentence that connects Part III back
to Parts I and II.

\emph{Cuba's welfare state has never been financed by Cubans.}

It was financed between 1972 and 1989 by Soviet terms of trade --- sugar bought
above the world price, oil supplied below it, and a petroleum rent Cuba was
permitted to re-export. It was financed from 2004 to about 2015 by Venezuelan
oil and by the medical-services contracts that oil relationship generated. In
the intervals it was financed by monetary emission and by not maintaining
anything, which is why the hospitals have no water and Havana loses buildings to
collapse.

An entitlement funded by an external rent lasts exactly as long as the rent. Cuba
has now discovered this twice, in 1990 and again after 2015, and on both
occasions the welfare state was saved only by extraordinary effort and at the
cost of its physical plant. The European systems this book takes as its standard
are financed by taxing a domestic base --- which is why they survive changes of
government, of commodity price, and of patron, and why no European government
has ever had to choose between paying pensions and importing fuel.

So the proposal is: same benefit, different plumbing. A European benefit
structure on a European tax structure. Rent-financed becomes tax-financed.

This is not a smaller undertaking than it sounds. Cuba has no modern tax system
--- it abolished most taxation in the 1960s on the grounds that a state which
owns everything need not tax itself, and reintroduced taxes only from 1994 and
only partially. It has no culture of filing, no significant enforcement history,
no withholding infrastructure outside the state payroll, and a population that
has never in living memory paid income tax. Building a revenue administration
that can collect a value-added tax has taken five to eight years everywhere it
has been attempted. That is the single longest lead time in this book and it is
why Chapter~\ref{ch:sequence} starts it in the first phase, when it yields
nothing.

\section{The arithmetic}

The numbers below are orders of magnitude with their assumptions stated, not
forecasts. They are here because a commitment of this kind is meaningless
without a rough answer to what it costs and whether it can be paid for.

\emph{What the benefit costs.} A European-standard universal system runs at
roughly 25 to 30 per cent of national income in social spending: public health
around 7 to 8 per cent, education around 5, pensions around 10 to 12, other
social protection the balance. Total general government spending in EU member
states sits in the 40 to 50 per cent range and tax revenue in the 35 to 45 per
cent range.\unv{}

\emph{What Cuba currently does.} Cuban published figures put social spending at
a very high share of output and revenue likewise, but --- per the note on
sources --- most of that is an artefact of converting a state-owned economy at
an administrative exchange rate, and the ratios do not mean what the same ratios
mean elsewhere. What can be said without the conversion problem is physical:
Cuba currently delivers something close to universal coverage at a real resource
cost far below European levels, because it pays its doctors, nurses and teachers
a fraction of a European wage and has stopped maintaining its buildings and
replacing its equipment. The system is being subsidised by its own staff and its
own capital stock.

\emph{Where the gap is.} That is the honest form of the problem. Cuba does not
need to buy European coverage --- it has coverage. It needs to buy European
\emph{inputs}: wages that stop the staff leaving, drugs that are actually in the
pharmacy, equipment that works, buildings that do not flood, and a pension that
someone can live on. Those are the expensive parts and they are exactly the
parts currently unfunded.

\emph{Whether it is affordable now.} No. At Cuban 2026 output a European benefit
level is not financeable at any tax rate, and a book that pretended otherwise
would be worthless. What is proposed instead is a stated trajectory: protect the
\emph{floor} immediately and let the \emph{ceiling} converge as the base grows,
with the convergence given a target and a date rather than left aspirational,
and with the tax system built now so that it can collect the revenue when the
base exists. The design below is written to be affordable in its first phase and
to scale.

\section{Which European model}

Three families, and the choice matters more than it looks.

The \emph{Nordic} model is universal, financed largely from general revenue with
a high broad-based consumption tax, service-heavy rather than cash-heavy, and
depends on a high employment rate. The \emph{Bismarckian} model of Germany and
France is contributory and payroll-financed, so coverage follows formal
employment. The \emph{Southern European} pattern is universal in statute and
patchy in delivery, financed by a narrow base alongside large informality.

Cuba would drift into the third by default. It is therefore the outcome to
design against, and the design choices below are made specifically to avoid it.

The book proposes a \emph{Nordic benefit structure} --- universal, service-heavy,
not means-tested --- financed by a \emph{broad consumption and property base}
rather than by payroll.

Universality rather than means-testing, for two reasons that are about Cuba
rather than about principle. Cuba already has universality and dismantling it to
build a targeting apparatus would destroy a working system to construct one it
has no capacity to run. And every means test is a discretionary decision taken
by an official about a citizen, which is precisely the surface
Chapter~\ref{ch:integrity} spends its length trying to eliminate. In a country
where \emph{resolver} is a verb, a benefit that depends on an official's
assessment is a benefit that depends on knowing the official.

Against payroll financing as the main instrument, for two more. Coverage would
track formal employment in an economy where most new jobs will start informal,
which reproduces the Southern European outcome directly. And a high payroll
wedge is the most reliable known way to kill small enterprise, which is the
sector Chapter~\ref{ch:capital} needs to grow and which Cuba has spent fifty
years failing to have.

\begin{design}{The benefit structure}
\label{d:benefit}
\begin{rules}
\rl{1}{Universal in kind} Health care and education free at the point of use to
every resident, without means test, contribution condition or nationality
qualification beyond residence. This is the existing Cuban settlement and it is
retained unchanged.

\rl{2}{Universal in cash} A universal flat basic pension at a level set to
prevent poverty, payable to every resident above the retirement age regardless
of contribution record; a universal child benefit; and disability support on
assessed need rather than on contribution.

\rl{3}{Earnings replacement separately} Above the flat tier, an earnings-related
pension on individual accounts, funded or notional, with the contribution
visible on the payslip. Separating the two tiers is what allows the flat tier to
be adequate without being unaffordable.

\rl{4}{Unemployment} Time-limited unemployment insurance with an active
component, introduced in the phase in which state employment is reduced and not
before --- because it is the instrument that makes that reduction politically
possible, and introducing it late is what turned the 2011 payroll plan into a
plan nobody executed.

\rl{5}{Legislated, not administered} Every benefit level is set by legislation
and uprated by a published formula. No benefit is variable at administrative
discretion. This is the difference between an entitlement and a favour, and per
Design~\ref{d:commit} it is also a commitment device.
\end{rules}
\end{design}

\section{The tax system}

This is the operational core of the chapter. Chapter~\ref{ch:money} owns
stabilisation, the currency and the deficit; this chapter owns what is taxed and
at what rate. The two must agree on one number --- the revenue share the welfare
commitment requires --- and Chapter~\ref{ch:money} is where it is shown to be
financeable without inflation.

The design principle throughout is that a weak administration collects a simple
tax and fails at a clever one. Cuba is building a revenue service from almost
nothing, in a cash economy, with a demoralised civil service and no filing
culture. Every sophistication is a place where the system leaks or where an
official acquires discretion.

\begin{design}{The revenue structure}
\label{d:tax}
\begin{rules}
\rl{1}{Broad base, low rates, no exemptions} Stated first because it governs
everything below. Exemptions are how tax systems die: each is small, each is
defensible, each creates a boundary to be argued over, and the argument is
conducted between a taxpayer and an official with discretion. No sectoral
holidays, no negotiated rates, no special regimes for named investors.

\rl{2}{VAT as the workhorse} A single-rate value-added tax, one rate and no
reduced rates, with a registration threshold high enough that micro-enterprise
stays outside it. Food is \emph{not} zero-rated: zero-rating is regressive per
peso of revenue forgone --- the rich buy more food than the poor in absolute
terms --- and the distributional repair belongs on the spending side, through
clause~2 of Design~\ref{d:benefit}, where it can be aimed.

\rl{3}{Income tax, withheld and simple} Progressive, few brackets, a high exempt
threshold so that most workers file nothing, and withholding at source wherever
a source exists. The objective in the first decade is coverage and habit, not
yield.

\rl{4}{Property tax} An annual tax on residential and commercial property,
assessed on area and location by formula rather than by valuation judgment,
collected locally and retained locally per Design~\ref{d:municipal}. This is the
easiest tax to administer, the hardest to evade, and the natural local revenue.
It is also the instrument that makes the 1960 Urban Reform's gift pay for the
services around it: Cuba has near-universal home ownership and a collapsing
housing stock, and the two facts are the same fact.

\rl{5}{Presumptive regime for micro-enterprise} A single small payment based on
turnover, replacing income tax, VAT and social contributions together, on the
Latin American \emph{monotributo} model, with a published graduation path into
the ordinary regime. Get this wrong and the private sector stays informal
permanently, which is the Southern European failure this chapter is designed
against.

\rl{6}{Mobile bases taxed lightly and simply} Remote-work and nomad income under
Chapter~\ref{ch:talent} pays a low flat rate with no deductions. A high rate on
a base that can leave collects nothing; a low rate on a base that stays collects
something and is worth complying with.

\rl{7}{Rents where they exist} An accommodation levy on visitor nights, lease
and concession revenue from state land and coastline, and nickel royalties. See
Design~\ref{d:rentrule} for what may be done with the proceeds, which is the
part that matters.

\rl{8}{What is not done} No wealth tax --- unadministrable here and it drives
out exactly the diaspora capital Chapter~\ref{ch:capital} is trying to attract.
No exit tax on emigrants, which is a tax on the grievance this book is trying to
settle. No high statutory corporate rate with negotiated exemptions, which is a
corruption machine with a tax attached. No earmarking beyond
Design~\ref{d:floor}, because earmarking fragments a budget and then binds it in
the wrong year.
\end{rules}
\end{design}

\section{The rent trap}

Parts I and II show the same failure three times: the American sugar quota, the
Soviet protocols, the Venezuelan oil. A fourth is available --- a tourism boom,
a normalisation windfall, a nickel price spike --- and if it arrives it will be
spent on recurrent social commitments, because that is what every government in
this position has done, including this one.

The welfare state must be financed from a broad domestic base \emph{because} a
rent-financed one has already died twice in living memory. That is not a
preference for taxation on ideological grounds. It is the specific institutional
lesson of Chapters~\ref{ch:sovieteconomy} and \ref{ch:venezuela}.

\begin{design}{The rent rule}
\label{d:rentrule}
\begin{rules}
\rl{1}{Separation} Revenue from concessions, royalties, privatisation proceeds,
debt relief and any windfall is paid into a stabilisation and investment fund
and is not available to the recurrent budget.

\rl{2}{Permitted uses} The fund finances capital spending, the pension
transition of Design~\ref{d:pension}, and countercyclical support after a
hurricane or an external shock. It does not finance wages, benefits or
subsidies.

\rl{3}{Anchor} The recurrent budget balances over the cycle on non-rent revenue
alone, with the deficit rule published, monitored by an independent fiscal
council, and reported on annually to the assembly.

\rl{4}{Escape} The rule may be suspended only for a declared natural disaster or
external shock, by supermajority, with an automatic expiry and a published
return path. A fiscal rule with a discretionary escape hatch is not a rule, and
a fiscal rule with no escape hatch at all does not survive a Category 5
hurricane --- which Cuba will have.
\end{rules}
\end{design}

\section{The pension arithmetic}

This is the hardest number in the book and the chapter should not pretend
otherwise.

Cuban fertility has been below replacement since 1978, is now around 1.2 to 1.4,
and there is no cohort in the pipeline. More than a quarter of the population is
over sixty. And between one and two million people, disproportionately of
working age, left between 2021 and 2025. A pay-as-you-go pension at European
replacement rates is not financeable on that dependency ratio at any plausible
tax rate, and no amount of design changes that.

What design can do is decide who bears the gap and make the arithmetic visible
in advance rather than in a crisis.

\begin{design}{Pensions}
\label{d:pension}
\begin{rules}
\rl{1}{Two tiers} The universal flat pension of Design~\ref{d:benefit},
tax-financed, set to prevent poverty and uprated by formula; plus an
earnings-related tier on individual accounts. The flat tier is what protects the
people Chapter~\ref{ch:premises} identified as exposed. The earnings tier is
what stops the flat tier from having to be generous.

\rl{2}{Retirement age indexed} The pension age is indexed to life expectancy by
a formula legislated once and applied automatically thereafter. Legislating the
formula is politically possible; legislating each increase separately is not,
which is why countries that chose the second route do not have the increases.

\rl{3}{Immigration is pension policy} The cheapest instrument for a dependency
ratio is more working-age people. Return migration and immigration are therefore
treated in Chapter~\ref{ch:talent} as fiscal instruments and not as a cultural
question, and the residence regime there is designed with this chapter's
arithmetic in view.

\rl{4}{Portability} Bilateral social-security agreements with the countries
holding the diaspora --- the United States, Spain, Mexico, Italy --- so that
contributions made abroad count, returnees are not penalised, and Cubans abroad
can contribute voluntarily. This converts a diaspora into a contributor base,
which no other instrument does.

\rl{5}{Honest accounting} An independent actuarial report to the assembly every
year, published, stating the projected gap and what closes it. Some of this gap
does not close, and a published number is how a country decides which part it
will meet by later retirement, which by lower replacement rates, which by
higher taxes and which by migration. Cuba's pension deficit is currently
neither published nor projected.
\end{rules}
\end{design}

\section{The staff}

A welfare state is people, and this is where the current arrangement is failing
fastest. Cuban doctors, nurses and teachers left for the tourism sector in the
1990s and are leaving the country in the 2020s. No financing plan works if the
staff are gone, and no plan that treats staffing as free has been costed.

\begin{design}{Public pay}
\label{d:pay}
\begin{rules}
\rl{1}{Benchmark} Public professional pay is set against the domestic private
and tourism wage, published, and reviewed annually by an independent body. It is
not set against a historic peso scale, which is how a surgeon came to earn less
than a bellhop and how the whole society learned that its education system was
a joke.

\rl{2}{Consistency} Chapter~\ref{ch:tourism} and Chapter~\ref{ch:life} are bound
by this clause. A tourism strategy that bids up service wages without a
corresponding public pay settlement is a strategy for emptying the hospitals,
and Cuba has run that experiment.

\rl{3}{Costed} The public wage bill under clause~1 is the largest single line in
this chapter's cost and is stated as such in Chapter~\ref{ch:sequence}'s
financing table rather than treated as an administrative detail.

\rl{4}{Retention before recruitment} In the first phase, spending goes to
retaining existing staff rather than to training replacements, because a
trained nurse who leaves takes four years to replace and the outflow compounds.
\end{rules}
\end{design}

\section{Delivery, and where the line falls}

Universal coverage does not require state monopoly provision, and pretending it
does costs Cuba capacity it cannot spare. But a two-tier system in the wrong
place destroys the thing being protected, because once the professional and
political classes leave the public system it stops being defended in the budget.
So the line has to be drawn explicitly rather than left to practice.

\begin{design}{Public and private provision}
\label{d:delivery}
\begin{rules}
\rl{1}{Private provision permitted} In dentistry, pharmacy, optical services,
elective and cosmetic procedures, eldercare, rehabilitation, and non-compulsory
tertiary and vocational education. In these, private and cooperative providers
operate alongside the public system, are licensed and inspected, and may
contract to it.

\rl{2}{Public provision protected} Primary care, emergency care, maternity,
public health and immunisation, and compulsory schooling remain universal,
public and free at the point of use, with no parallel private tier permitted to
substitute for them. These are the services whose universality produces the
outcomes in Chapter~\ref{ch:ledgerch}, and they are the ones a two-tier
settlement hollows out first.

\rl{3}{No queue-jumping} No public facility sells priority access. Health
tourism under Chapter~\ref{ch:life} operates in designated facilities with
ring-fenced staff, and its revenue accrues to the public system --- and if that
ring-fence cannot be maintained, the programme stops, because
Chapter~\ref{ch:premises} established that staff are the binding scarcity.

\rl{4}{Cooperatives first} Where provision is opened, cooperative and
professional-partnership forms are the default and are given the licensing
preference, because they distribute the surplus to the practitioners and are
the form least likely to concentrate.
\end{rules}
\end{design}

\section{Making it stick}

Everything above is a policy until it is entrenched, and Chapter~\ref{ch:polity}
established what happens to Cuban policies.

\begin{design}{The floor}
\label{d:floor}
\begin{rules}
\rl{1}{Constitutional minimum} A floor on combined health and education spending
as a share of general government expenditure, set in the constitution, below
which a budget may not be presented. This is the one earmark
Design~\ref{d:tax}.8 permits, and it is permitted because it is the commitment
device the whole chapter rests on.

\rl{2}{Separated accounts} Social insurance accounts are held separately from
the general budget, published in full, and may not be borrowed from.

\rl{3}{Annual actuarial report} Per Design~\ref{d:pension}.5, to the assembly,
published, independent.

\rl{4}{Justiciable} The rights to health care and to compulsory schooling are
enforceable by an individual in the administrative court of
Design~\ref{d:courts}. Not the level of provision --- courts should not set
budgets --- but the fact of access, and the non-discrimination of it.

\rl{5}{Published outcomes} Health and education outcomes are published annually
by province and, per Chapter~\ref{ch:premises}, disaggregated by race. Cuba
declared the racial question solved in 1962 and then made it unmeasurable for
thirty years, and entered its market opening with no idea where it stood. The
first requirement of this design in that domain is to publish the data.
\end{rules}
\end{design}

\section{What this chapter costs the rest of the book}

Stated plainly, because a design that hides its trade-offs is a brochure.

A European welfare state means a tax burden well above the Latin American norm.
Cuba will be bidding for investment against countries in the same region with
lower rates, weaker labour protection and no comparable social spending, and it
will lose some of that bidding. Chapter~\ref{ch:capital} therefore cannot
compete on rates and must compete on the other things investors price:
certainty, speed of approval, enforceable contracts, an honest customs service,
a court that can be lost in fairly, and a workforce that can read. Those are
harder to build than a tax holiday and they are also, unlike a tax holiday,
worth having for their own sake.

It also means that the growth chapters have to work. A welfare state at this
standard requires a tax base that does not currently exist, which is to say it
requires Chapters~\ref{ch:energy} through \ref{ch:capital} to succeed. This
chapter is not an alternative to the growth strategy. It is the reason for it.
